% F035 指数幂的运算法则｜02 快懂：先明白原因
\begin{explanationbox}
    \textbf{\textcolor{CExplanation}{为什么要学这三条？}}\par
    算 $2^3\times2^4$，若先求出 $8\times16$，还要再乘一次；如果直接看指数，$3+4=7$，马上得到 $2^7$。\textbf{它把难算的连乘，变成容易算的加法或乘法。}

    \medskip
    \textbf{\textcolor{CExplanation}{一、 为什么同底相乘，要“加”？}}\par
    $2^3$ 是 3 个 $2$ 相乘，$2^4$ 是 4 个 $2$ 相乘。两串接在一起，一共有 $3+4=7$ 个 $2$。这里\textbf{不是“指数都要加”，而是相同底数的幂在相乘}。
    \[2^3\cdot2^4=\underbrace{2\cdot2\cdot2}_{3\text{ 个}}\cdot\underbrace{2\cdot2\cdot2\cdot2}_{4\text{ 个}}=2^7.\]

    \medskip
    \textbf{\textcolor{CExplanation}{二、 为什么幂的乘方，要“乘”？}}\par
    $(2^3)^4$ 是把 $2^3$ 连乘 4 次。\textbf{每组 3 个 $2$，一共 4 组}，所以共有 $3\times4=12$ 个 $2$。
    \begin{center}
    \begin{tikzpicture}[font=\small]
        \foreach \x in {0,2.25,4.5,6.75} {
          \node[draw=blue!50!black,fill=blue!6,rounded corners=3pt,minimum width=2.0cm,minimum height=0.7cm] at (\x,0) {$2\cdot2\cdot2$};
        }
        \node[font=\footnotesize] at (3.4,-0.66) {每组 3 个，共 4 组：$3\times4=12$ 个};
        \node[draw=green!50!black,fill=green!10,rounded corners=3pt] at (3.4,-1.4) {$(2^3)^4=2^{12}$};
    \end{tikzpicture}

    \end{center}
    \textbf{关键区别：} “两串连起来”数相加；“每组几个、总共几组”数相乘。

    \medskip
    \textbf{\textcolor{CExplanation}{三、 为什么积的乘方，可以分开？}}\par
    $(2\times3)^3=(2\times3)(2\times3)(2\times3)$。把三个 $2$ 放一起、三个 $3$ 放一起，结果不变：
    \[(2\cdot3)^3=\underbrace{2\cdot2\cdot2}_{2^3}\cdot\underbrace{3\cdot3\cdot3}_{3^3}=2^3\cdot3^3.\]
    \begin{center}
    \begin{tikzpicture}[font=\small,>=stealth]
        \node[draw=orange!70!black,fill=orange!7,rounded corners=3pt,inner sep=6pt] (a) at (0,0) {$(2\cdot3)(2\cdot3)(2\cdot3)$};
        \node[draw=green!50!black,fill=green!10,rounded corners=3pt,inner sep=6pt] (b) at (6,0) {$(2\cdot2\cdot2)(3\cdot3\cdot3)$};
        \draw[->,thick,green!50!black] (a.east)--node[above,font=\footnotesize]{交换位置}(b.west);
        \node at (3,-0.75) {$(2\cdot3)^3=2^3\cdot3^3$};
    \end{tikzpicture}

    \end{center}

    \textbf{\textcolor{CExplanation}{从整数到分数指数，为什么还能算？}}\par
    $a^0=1$（$a\ne0$），$a^{-1}=1/a$，$a^{1/2}=\sqrt a$（$a>0$）。这些新写法仍与原来的规律一致。例如：
    \[2^3\cdot2^{-1}=8\cdot\frac12=4=2^2,\qquad 9^{1/2}\cdot9^{1/2}=3\cdot3=9.\]
    当指数不是整数时，不能再说“重复乘了几次”；\textbf{先用整数例子理解来源，再按定义与成立条件使用实数指数运算法则}。

    \medskip
    \textbf{\textcolor{CExplanation}{考试中什么时候最快？}}\par
    看到 $a^m a^n$ 就想“指数加”；看到 $(a^m)^n$ 就想“指数乘”；看到 $(ab)^n$ 就想“两边都乘方”。底数不同，先看能否\textbf{化同底}；指数相同，想想能否\textbf{合并底数}。
\end{explanationbox}
